\section{Requirement conventions}

\guidance{Every requirement should be measurable or clearly verifiable.
Avoid ``should be robust'' without a number or pass/fail test.}

\subsection{ID scheme (by department)}

\noindent\begin{tabularx}{\linewidth}{@{}l X@{}}
\toprule
\textbf{Prefix} & \textbf{Department / scope} \\
\midrule
\texttt{MIS-} & Mission / campaign success \\
\texttt{VEH-} & Vehicle system-level (cross-department) \\
\texttt{ENG-} & Engine \\
\texttt{GSE-} & Ground Support Equipment \\
\texttt{REC-} & Recovery \\
\texttt{AERO-} & Aerostructure \\
\texttt{AV-} & Avionics \\
\texttt{SIM-} & Simulations \\
\texttt{SAF-} & Safety (cross-cutting) \\
\bottomrule
\end{tabularx}

\subsection{Priority}

\begin{itemize}
  \item \textbf{Must} --- required for mission success / safety; non-compliance blocks flight
  \item \textbf{Should} --- important; waiver needs lead approval
  \item \textbf{Could} --- desirable; may slip
\end{itemize}

\subsection{Status}

\begin{itemize}
  \item \textbf{Proposed} / \textbf{Approved} / \textbf{Frozen} / \textbf{Waived}
\end{itemize}

\subsection{V\&V methods}

The \textbf{V\&V} column in requirement tables is the planned
\emph{verification and validation} method (how we prove the requirement).
Codes may be combined (e.g.\ \texttt{T/A}).

\begin{table}[ht]
\centering
\caption{V\&V method codes used in requirement tables.}
\label{tab:vv-methods}
\begin{tabularx}{\linewidth}{@{}c l X@{}}
\toprule
\textbf{Code} & \textbf{Method} & \textbf{Meaning / example} \\
\midrule
T & Test & Measurement or experiment (flight, static fire, cold-flow, drop test) \\
A & Analysis & Calculation or simulation (trajectory, loads, margins) \\
I & Inspection & Visual / document check (post-flight damage, packing list, waiver on file) \\
D & Demonstration & Show the procedure works (abort drill, arming sequence, remote safing) \\
\bottomrule
\end{tabularx}
\end{table}

\subsection{Requirement table columns}

Each requirement table uses: \textbf{ID}, \textbf{Requirement}, \textbf{Priority},
\textbf{V\&V} (see Table~\ref{tab:vv-methods}), \textbf{Owner}.
